\begin{textsample}{POS Dim 4 – sp – Score 6.00 – sp/sp\_inf\_e\_ef\_intro\_1.txt}  \label{ex:f4_pos_221}
Introdução O Estado de São Paulo : números que impressionam!
O Estado de São Paulo é dividido atualmente em 645 municípios que ocupam um total de pouco mais de 248.000 km² de \textbf{território}, o que representa apenas 2,9 \% da superfície terrestre brasileira.
No entanto, os municípios paulistas contam, juntos, com aproximadamente 45 milhões de habitantes, de acordo com as estimativas populacionais do Instituto Brasileiro de Geografia Estatística ( IBGE ) de março de 2019.
Esse montante corresponde a mais de 22 \% do total da \textbf{população} do nosso país.
A \textbf{população} paulista é superior à de países como o Canadá ( 35.881.659 ), Peru ( 31.331.228 ) e Ucrânia ( 43.952.299 ) e muito próxima à da Argentina ( 44.694.198 ), de acordo com o Atlas The World Fact Book, da CIA — Central Intelligence Agency (2018)¹. ¹ Informações disponíveis em : [ https://www.cia.gov/library/publications/resources/the-world-factbook/fields/335rank.html](https://www.cia.gov/library/publications/resources/the-world-factbook/fields/335rank.html ).
A \textbf{população} paulista é uma das mais diversificadas e descende principalmente de africanos, \textbf{indígenas}, italianos, portugueses e de \textbf{migrantes} de outras regiões do país.
Outras grandes correntes imigratórias, como a de árabes, \textbf{alemães}, espanhóis, japoneses e chineses, tiveram presença significativa na composição étnica e cultural da \textbf{população} do estado.
Esses dados mostram quão diversa é a \textbf{população} paulista, assim como a enormidade do quantitativo de pessoas que habitam um espaço tão pequeno, quando comparado às dimensões continentais do país.
Na Educação Básica, as matrículas nas diferentes redes atingem o total de 7.433.331, segundo dados coletados no Cadastro de Alunos em fevereiro de 2019.
As tabelas 1 a 5 apresentam a distribuição dessas matrículas, permitindo mais acurada avaliação da dimensão das redes e da quantidade de crianças e estudantes atendidos.
Tabela 1 — Distribuição das crianças e estudantes matriculados na Educação Básica | Rede | Matrículas | |---|---:| | Privada | 699.954 | | Estadual | 3.241.473 | | Municipal | 3.491.994 | | Total | 7.433.421 | Fonte : Cadastro de Alunos, fev. 2019 Tabela 2 — Distribuição das crianças matriculadas na Educação Infantil | Rede | Matrículas | |---|---:| | Particular | 324.072 | | Estadual | 69 | | Municipal | 1.279.461 | | Total | 1.603.602 | Fonte : Cadastro de Alunos, fev. 2019 Tabela 3 — Distribuição dos estudantes matriculados no Ensino Fundamental — Anos Iniciais | Rede | Matrículas | |---|---:| | Particular | 91.068 | | Estadual | 646.725 | | Municipal | 1.667.015 | | Total | 2.404.808 | Fonte : Cadastro de Alunos, fev. 2019 Tabela 4 — Distribuição dos estudantes matriculados no Ensino Fundamental — Anos Finais | Rede | Matrículas | |---|---:| | Particular | 60.150 | | Estadual | 1.390.583 | | Municipal | 532.619 | | Total | 1.983.352 | Tabela 5 — Distribuição dos estudantes matriculados no Ensino Médio | Rede | Matrículas | |---|---:| | Particular | 224.664 | | Estadual | 1.204.096 | | Municipal | 12.899 | | Total | 1.441.569 | Fonte : Cadastro de Alunos, fev. 2019 Cabe destacar que os números representados nas tabelas anteriores não incluem aos estudantes matriculados na Educação de Jovens e Adultos e na Educação Profissional.
Somadas essas matrículas, segundo o Censo Escolar de 2015, o total de estudantes matriculados em escolas no Estado de São Paulo chega ao número 10.051.652.
Tais dados ressaltam, por um lado, o desafio enfrentado pelo Estado de São Paulo para assegurar educação de qualidade a todos os estudantes matriculados nas escolas paulistas ; por outro, a importância que o Currículo Paulista representa para a garantia de um patamar comum de aprendizagens.

% matched lemmas: alemão, indígena, migrante, população, território
\end{textsample}
